\chapter{Metodología}
\label{ch:metodologia}

\section{Datos y particiones}

El corpus procesado reside en \texttt{data/processed/} con particiones
\texttt{train.csv}, \texttt{val.csv} y \texttt{test.csv}. La release
\textbf{v6} (CDE \textbf{v1.9.0}) usa train $n=21{,}585$, val $n=459$ y test
$n=460$. El archivo \texttt{test.csv} conserva el mismo hash SHA-256 que en v1 y
B2, lo que garantiza comparabilidad directa entre checkpoints.

Cada instancia incluye el texto de entrada del modelo
(\texttt{texto\_modelo}), la etiqueta ordinal, metadatos de plataforma e
indicadores de sarcasmo, ironía, contexto político y conteo de jerga local
(\texttt{n\_jerga}).

\subsection{Evolución del corpus}

El checkpoint v1 entrenó con $n=2{,}144$ reales; B2 amplió el train a
$3{,}293$ instancias (sintéticos dirigidos). v6 integra anotación gold y
política \texttt{b2\_full} sobre el release congelado
\texttt{versiones/v6/dataset/}, inicializando pesos desde v4
(\texttt{versiones/v4/model}).

\subsection{Motivación histórica del reentrenamiento}

La auditoría del checkpoint v1 identificó cuatro huecos: atajo
emoji$\Rightarrow$tóxico, emojis de ataque no vistos, amenazas ofuscadas y sesgo
dialectal. Los releases B2 y v6 abordan robustez adversarial y equidad sin
cambiar el test oficial.

\section{Pipeline experimental}

El flujo de trabajo es reproducible y se alinea con los notebooks 01--07:

\begin{enumerate}[leftmargin=*]
  \item Análisis exploratorio (distribuciones, plataformas, longitudes).
  \item Normalización de texto y extracción de features (pesos de clase,
        indicadores dialectales).
  \item Baselines TF-IDF (Logistic Regression, LinearSVC, MultinomialNB).
  \item Fine-tuning de mBERT~\cite{devlin2019bert} sobre
        \texttt{versiones/v6/model} (checkpoint publicado en Hugging Face
        \texttt{caeher/mbert-sv}).
  \item Evaluación completa: métricas globales, por clase, por cortes, SLO de
        producción y comparación descriptiva de modelos.
  \item Benchmark zero-shot sobre datasets públicos de Hugging Face (checkpoint
        v1; no re-ejecutado en v6).
  \item Interpretabilidad con LIME~\cite{ribeiro2016lime} y
        SHAP~\cite{lundberg2017shap}.
\end{enumerate}

\section{Métricas de evaluación}

Se implementa un conjunto amplio de métricas en el módulo de evaluación del
proyecto:

\begin{itemize}[leftmargin=*]
  \item \textbf{Clásicas:} accuracy, balanced accuracy, precision/recall/F1
        (macro y weighted), MCC.
  \item \textbf{Ordinales:} MAE, RMSE y porcentaje de errores grandes
        ($\lvert y - \hat{y}\rvert \ge 2$).
  \item \textbf{Acuerdo:} kappa lineal y Quadratic Weighted Kappa (QWK).
  \item \textbf{Probabilísticas:} ROC-AUC OVR macro, log-loss, ECE y Brier macro.
  \item \textbf{Incertidumbre:} intervalo de confianza bootstrap del F1-macro
        (1000 remuestreos, $\alpha = 0{,}05$).
  \item \textbf{Score operativo:} agregado definido en
        \texttt{production\_slo\_v1.json} (cortes mínimos, recall de amenazas y
        restricciones adversariales).
\end{itemize}

La métrica de decisión principal es el \textbf{F1-macro}, que trata por igual
las cuatro clases y es adecuada ante desbalance.

\subsection{Pisos triviales sin modelo}

Los baselines de referencia (clase mayoritaria, azar uniforme y azar
estratificado) se calculan sobre la distribución del \textbf{test} ($n=460$) con
$k=4$ clases fijas.

\section{Cortes (slices)}

Además del desempeño global, se calculan métricas sobre subconjuntos
definidos en el test (plataforma, jerga, sarcasmo, ironía y contexto político).
Los cortes de mBERT v6 se derivan del vector de predicciones oficial
(\texttt{mbert\_test\_predictions.csv}, checkpoint \texttt{checkpoint-4050}) y
se exportan en \texttt{reports/metrics/slice\_metrics.csv}.

\subsection{Batería adversarial diagnóstica}

Existen dos conjuntos adversariales: \texttt{test\_adversarial.csv} ($n=75$) y
\texttt{test\_adversarial\_144.csv} ($n=144$, familias A--G). Son
\textbf{exclusivamente diagnósticos}: no entrenan ni seleccionan el modelo, y no
sustituyen el F1-macro oficial sobre \texttt{test.csv}.

\section{Contrato de inferencia}

Entrenamiento, evaluación por lote e inferencia del prototipo XLM-R comparten un contrato versionado (\texttt{inference\_contract.json}): \texttt{max\_length}=128, columna \texttt{texto\_modelo} y normalizador \texttt{normalize\_for\_model} v1.1. La auditoría reportó F1-macro 0.8058 con evaluación a \texttt{max\_length}=256; bajo el contrato oficial (\texttt{max\_length}=128, normalizador normalize\_for\_model v1.1), el F1-macro es 0.8058 (0.0\% de predicciones distintas frente a la ruta diagnóstica a 256, F1=0.8058). Sin normalización sobre texto crudo, 1.74\% de las predicciones cambian respecto al contrato.

